\chapter{焦点标的与验证路径}

本章聚焦能够改变下一轮配置排序的公司，而不是重述142只候选。我们把标的分为可复核当前价模型、恢复性估值模型和高空间但证据不足的验证观察三类；任何一类都必须有可观察的升级条件和明确的反证。

\section{可复核当前价模型}
\subsection{恒逸石化（000703）｜可复核正向验证模型}
现价14.48元，本机构概率值23.42元，对应61.7\%空间。熊/基准/牛值为10.86/25.59/40.41元；外部锚17.05元，权重10\%。
我们将其定义为可复核正向验证模型而非一般观察名单，原因是当前价格、股本、情景分母和外部证据可以同时复核。下一轮的关键验证是产品价差、原料成本、装置运行率、库存和现金流；若出现价差收窄、原料成本上升、装置检修或盈利恢复无法持续，则该模型应被降级。

\subsection{宏桥控股（002379）｜可复核正向验证模型}
现价18.66元，本机构概率值25.80元，对应38.3\%空间。熊/基准/牛值为14.00/26.24/40.65元；外部锚29.00元，权重10\%。
我们将其定义为可复核正向验证模型而非一般观察名单，原因是当前价格、股本、情景分母和外部证据可以同时复核。下一轮的关键验证是金属价格、产量、单位成本、库存和经营现金流；若出现价格回落、成本曲线抬升、产量不达预期或利润无法转化为现金，则该模型应被降级。

\subsection{亿纬锂能（300014）｜可复核正向验证模型}
现价52.28元，本机构概率值60.43元，对应15.6\%空间。熊/基准/牛值为31.49/68.46/74.57元；外部锚77.00元，权重10\%。
我们将其定义为可复核正向验证模型而非一般观察名单，原因是当前价格、股本、情景分母和外部证据可以同时复核。下一轮的关键验证是出货量、订单、产品价格、产能利用率和营运资金；若出现订单转化慢、价格竞争、利用率不足或库存/应收继续占用现金，则该模型应被降级。

\subsection{川能动力（000155）｜下行纪律对照}
现价11.09元，本机构概率值10.04元，对应-9.5\%空间。熊/基准/牛值为5.43/9.03/13.67元；外部锚20.43元，权重10\%。
我们将其定义为下行纪律对照而非一般观察名单，原因是当前价格、股本、情景分母和外部证据可以同时复核。下一轮的关键验证是利用小时、上网电价、燃料/资源成本、资本开支和现金流；若出现利用率或电价下行、成本抬升、资本开支超预期或现金流不达标，则该模型应被降级。

\noindent
\begin{exhibitbox}[图表 8｜正式模型横向比较：价值空间必须接受同一组验证纪律]
\small
\begin{tabularx}{\textwidth}{L{2.1cm}L{2.4cm}X L{4.0cm}}
\toprule
\textbf{标的} & \textbf{现价/概率值} & \textbf{本轮定位} & \textbf{下一轮必须验证} \\
\midrule
恒逸石化（000703） & 14.48/23.42元 & 可复核正向验证模型；61.7\%空间 & 验证：产品价差、原料成本、装置运行率、库存和现金流；反证：价差收窄、原料成本上升、装置检修或盈利恢复无法持续 \\
宏桥控股（002379） & 18.66/25.80元 & 可复核正向验证模型；38.3\%空间 & 验证：金属价格、产量、单位成本、库存和经营现金流；反证：价格回落、成本曲线抬升、产量不达预期或利润无法转化为现金 \\
亿纬锂能（300014） & 52.28/60.43元 & 可复核正向验证模型；15.6\%空间 & 验证：出货量、订单、产品价格、产能利用率和营运资金；反证：订单转化慢、价格竞争、利用率不足或库存/应收继续占用现金 \\
川能动力（000155） & 11.09/10.04元 & 下行纪律对照；-9.5\%空间 & 验证：利用小时、上网电价、燃料/资源成本、资本开支和现金流；反证：利用率或电价下行、成本抬升、资本开支超预期或现金流不达标 \\
\bottomrule
\end{tabularx}
\normalsize
\sourcenote{正式模型均要求当前价格、分母、原始外部锚和失效条件同时可复核；若概率值低于现价，则仅承担下行纪律对照功能。}
\end{exhibitbox}

\section{恢复性估值：关注分母修复，而非追逐表面空间}

恢复模型的共同纪律是：当H1利润由低基数、扭亏、结算或周期集中驱动时，不允许直接以H1年化EPS乘常规行业PE。下表展示当前最具代表性的恢复样本及其下一步证据要求。

\noindent
\begin{exhibitbox}[图表 9｜恢复性估值：分母切换与验证条件]
\scriptsize
\begin{tabularx}{\textwidth}{L{1.1cm}L{1.7cm}L{2.4cm}R{1.0cm}R{1.0cm}X}
\toprule
\textbf{代码} & \textbf{标的} & \textbf{替代分母/方法} & \textbf{概率值} & \textbf{空间} & \textbf{升级条件与反证} \\
\midrule
601360 & 三六零 & 2026E 收入 PS；2026E收入95.91亿元与AI/安全变现 & 13.10 & 45.6\% & 验证：2026年收入接近95.91亿元，AI/安全形成可重复变现，EPS共识不低于0.06元且经营现金流转正；反证：AI产品变现不及预期、2026年EPS低于0.03元、经营现金流仍为负，或PS低于可比区间 \\
000042 & 中洲控股 & 项目结算正常化 EPS / PE；项目结算、税费与利息正常化EPS & 8.54 & 6.5\% & 验证：2026年EPS不低于1.10元，项目结算及税费/利息计提得到确认，债务与现金流压力改善且股价守住8元；反证：结算利润回转、土地增值税或利息计提上升、销售回款走弱，或净负债/流动性恶化 \\
002460 & 赣锋锂业 & 周期调整 EPS / PE；2026E周期调整EPS & 53.68 & 4.2\% & 验证：锂价、产量、资源自给率和H2经营现金流共同验证2026年EPS 4.45元；反证：锂价回落、资源项目延期、库存减值，或现金流弱于利润 \\
002812 & 恩捷股份 & 产能周期正常化前瞻 EPS / PE；隔膜价格/利用率恢复后的2027E EPS & 85.29 & 55.2\% & 验证：隔膜涨价、产能利用率、海外/3C订单和H2现金流共同验证2027年EPS 5.16元；反证：供需修复不及预期、价格回落、客户集中，或产能利用率不足 \\
300390 & 天华新能 & 周期调整 EPS / PE；锂价敏感的2026E周期调整EPS & 77.30 & 28.8\% & 验证：氢氧化锂价格、客户采购、产量和H2经营现金流共同验证2026年EPS 5.70元；反证：锂价回落、客户集中风险暴露，或经营现金流继续弱于利润 \\
002240 & 盛新锂能 & 周期调整 EPS / PE；锂价、自有矿与现金流桥接EPS & 37.89 & 21.8\% & 验证：锂价、出货、自有矿贡献和H2经营现金流共同验证2026年EPS 2.45元；反证：锂价回落、现金流持续为负，或资源项目贡献不及预期 \\
600736 & 苏州高新 & 正常化 EPS / PE（PB交叉检验）；园区/结算/投资收益正常化EPS & 7.19 & 4.7\% & 验证：产业园利润、地产结算、投资收益和经营现金流共同验证2026年EPS 0.14元；反证：地产结算延期、投资收益回落，或经营现金流继续为负 \\
000415 & 渤海租赁 & 正常化 EPS / PE（PB交叉检验）；租赁收入、资产质量与汇率正常化EPS & 4.40 & -1.1\% & 验证：租赁收入、资产质量、汇率和H2现金流共同验证2026年EPS 0.77元；反证：飞机减值、汇率损失、融资成本上升，或现金流风险上行 \\
\bottomrule
\end{tabularx}
\normalsize
\sourcenote{逐票正常化分母、可比集、倍数和公式代入见附录B；表中区间均为条件性本机构价值。}
\end{exhibitbox}

\section{高空间模型的正确使用方式}

高空间不是排序第一的理由。本轮已对全部6只100\%+空间样本完成逐票证据闭环：原始API、已归档PDF、目标价字段、Q1经营现金流和池内准入均已核验，闭环率为6/6。核验结果不是“仍待补资料”，而是当前正权原始目标为0只、正式升级为0只。

已找到的历史目标仅作为反向证据：九安医疗70.28元、吉林敖东约18.02元、辽宁成大22.79元、中山公用10.50--11.11元、中国船舶31.36元，均显著低于或仅接近内部高空间情景值，且因时效不足权重为0。江波龙在原始API及已归档PDF语料中未找到目标价；中国船舶50元媒体转载仅作零权重交叉验证。

\noindent
\begin{exhibitbox}[图表 9A｜高空间样本的已核验证据与准入结论]
\scriptsize
\begin{tabularx}{\textwidth}{L{0.95cm}L{1.45cm}R{0.95cm}L{5.0cm}X}
\toprule
\textbf{代码} & \textbf{标的} & \textbf{空间} & \textbf{已核验证据} & \textbf{准入结论/剩余事件} \\
\midrule
002432 & 九安医疗 & 323.4\% & API/PDF 2/2；当前原始目标0；2024-12-23 70.28元，权重0\%；Q1现金流-0.92亿元 & 保留候选观察，维持未进入优先池及非正式模型；未来：重新满足优先池的价格位置与行业阶段规则；H2扣非利润持续性与经营现金流转正 \\
000623 & 吉林敖东 & 317.0\% & API/PDF 2/2；当前原始目标0；2024-03-27 18.02元，权重0\%；Q1现金流0.11亿元 & 保留优先池验证，维持非正式模型；未来：2026-04-01之后原始券商目标价或合理价值区间；H2扣非利润持续性与投资收益质量 \\
600739 & 辽宁成大 & 300.0\% & API/PDF 4/4；当前原始目标0；2017-10-30 22.79元，权重0\%；Q1现金流-1.76亿元 & 保留优先池验证，维持非正式模型；未来：2026-04-01之后原始券商目标价或合理价值区间；H2扣非利润持续性与广发证券投资收益质量 \\
000685 & 中山公用 & 162.8\% & API/PDF 6/6；当前原始目标0；2025-03-04 10.50--11.11元，权重0\%；Q1现金流-0.27亿元 & 保留候选观察，维持未进入优先池及非正式模型；未来：重新满足优先池的价格位置与行业阶段规则；H2扣非利润与投资收益可持续性 \\
600150 & 中国船舶 & 142.9\% & API/PDF 56/8；当前原始目标0；2023-03-27 31.36元，权重0\%；Q1现金流81.52亿元 & 保留优先池验证，维持非正式模型；未来：当前原始券商目标价或合理价值区间；H2订单交付、船价、利润率与应收回款 \\
301308 & 江波龙 & 119.2\% & API/PDF 12/8；当前原始目标0；原始语料未找到可接受目标价；Q1现金流-28.75亿元 & 保留候选观察，维持未进入优先池及非正式模型；未来：重新满足优先池的价格位置与行业阶段规则；当前原始券商目标价或合理价值区间 \\
\bottomrule
\end{tabularx}
\normalsize
\sourcenote{专项证据包：\texttt{data/high\_upside\_evidence\_closure\_20260716.json}。当前原始目标缺失已按“核验未披露”关闭；H2利润、现金流、订单与价格周期保留为未来事件验证。}
\end{exhibitbox}
